The evaluation is confined to the DocVQA-2026 validation set (25 documents, 80 questions); the test set is held out by the competition portal, and the small validation size limits the resolution of per-category and per-slice estimates. The per-category and page-bucket analyses are computed on cross-model matched-configuration runs rather than the homogeneous baselines, so those figures are read as deltas, with the headline solver numbers ($n{=}8$) cited for absolutes. The CodeAct cells marked $\dag$ (the 8B, 9B, and Gemma rows) are produced by an earlier implementation of that scaffold and are provisional pending a re-run, whereas the 4B and 27B CodeAct cells are the corrected, final $n{=}8$ loop; the active-perception conclusions rest on the RLM numbers regardless.

Generalization is demonstrated across two model families. The training study is deliberately preliminary and narrowed in scope: it covers only the supervised stage, its quantitative validation scores are being re-run after an evaluation-harness error and so are not reported here, and the reinforcement learning and on-policy distillation set out in the proposal are left to future work. The account in \S\ref{sec:scaffold} of why a context-managing scaffold has a higher but currently inaccessible ceiling is a hypothesis rather than a demonstrated result.
